%%% FIELD proof-main
Put \(\nu=(d+1)/2\).  We first isolate an actual-design two-point subexperiment.  By the dimension, domain, fill, and separation assumptions, fix one admissible deterministic design sequence.  Fix a tuple \((c_X,c_W,\alpha_X,\alpha_W,\alpha_C)\in\mathcal K_{\mathrm{nuis}}\), as permitted by nuisance compactness, and fix \(p^\circ\in[p_{\min},p_{\max}]\), as permitted by the exponent-space assumption.  The Gaussian-field and spectral-family assumptions supply an admissible spectral matrix.  The determinant-margin condition and positivity of \(f_Xf_W\) then imply, for every admissible spectral matrix,
\[
 0<\rho\leq 1-\frac{f_{XW}(\omega)^2}{f_X(\omega)f_W(\omega)}\leq1.
 \tag{1}
\]
The amplitude-space assumption permits \(A=0\).  At this value \(f_{XW}=0\), so (1) gives
\[
 f_X(\omega)f_W(\omega)-f_{XW}(\omega)^2
 =f_X(\omega)f_W(\omega)
 \geq\rho f_X(\omega)f_W(\omega).
 \tag{2}
\]
Thus the determinant margin continues to hold.

For completeness, this choice gives a genuine Gaussian subexperiment.  Both marginal spectra are strictly positive and integrable: they are bounded near the origin, and their radial integrands at infinity have order
\[
 r^{d-1}r^{-(d+1)}=r^{-2}.
\]
For any finite set of sites and any real coefficient vector \(z\), the covariance form induced by either marginal spectrum \(f\in\{f_X,f_W\}\) is
\[
 \kappa_d\int_{\mathbb R^d}f(\omega)
 \left|\sum_i z_i e^{\mathrm i\omega^{\mathsf T}s_i}\right|^2\,d\omega\geq0,
 \tag{3}
\]
where \(\kappa_d>0\) is the common Fourier-normalization constant.  The resulting centered Gaussian finite-dimensional laws are consistent because restriction to fewer sites takes the corresponding principal covariance submatrix.  Taking zero cross-covariance gives independent centered stationary Gaussian fields \(X\) and \(W\) with the prescribed diagonal spectral matrix.  The observation-law assumption then gives \(Y=\beta X+W\).  Consequently, as \(\beta\) varies over \(\operatorname{int}(\mathcal B)\), these laws form an \(A=0\) subexperiment \(\mathcal E_0\subset\mathcal M\).  The value \(p^\circ\) is immaterial when \(A=0\), so this construction does not presume identification of \(p\), including at either exponent endpoint.

Let \(X_N=(X(s_{1,N}),\ldots,X(s_{N,N}))^{\mathsf T}\), define \(W_N\) analogously, and write their covariance matrices as \(\Sigma_{X,N}\) and \(\Sigma_{W,N}\).  The separation condition implies that the sites are distinct.  If \(z\ne0\), choose \(v\in\mathbb R^d\) outside the finite union of hyperplanes orthogonal to the site differences.  Then \(v^{\mathsf T}s_{i,N}\) are distinct.  If the exponential sum in (3) vanished identically, restricting to \(\omega=tv\) and differentiating at \(t=0\) for orders \(0,\ldots,N-1\) would give a nonsingular Vandermonde system and hence \(z=0\), a contradiction.  The sum is therefore nonzero on an open set.  Since both spectra are strictly positive, (3) is strictly positive for every \(z\ne0\).  Hence \(\Sigma_{X,N}\) and \(\Sigma_{W,N}\) have full rank, and the Gaussian laws below have full support on \(\mathbb R^{2N}\).

Define the finite positive constant
\[
 C_\star=\frac{c_X}{c_W}
 \max\left\{1,\left(\frac{\alpha_W^2}{\alpha_X^2}\right)^\nu\right\}.
 \tag{4}
\]
Its finiteness and positivity follow from nuisance compactness.  For every \(r\geq0\), the ratio \((\alpha_W^2+r)/(\alpha_X^2+r)\) lies between \(1\) and \(\alpha_W^2/\alpha_X^2\).  Therefore
\[
 \frac{f_X(\omega)}{f_W(\omega)}
 =\frac{c_X}{c_W}
 \left(\frac{\alpha_W^2+\|\omega\|^2}
 {\alpha_X^2+\|\omega\|^2}\right)^\nu
 \leq C_\star.
 \tag{5}
\]
Combining (3) and (5), for every \(z\in\mathbb R^N\),
\[
 z^{\mathsf T}\Sigma_{X,N}z
 \leq C_\star z^{\mathsf T}\Sigma_{W,N}z,
 \qquad\text{hence}\qquad
 \Sigma_{X,N}\preceq C_\star\Sigma_{W,N}.
 \tag{6}
\]
Let \(\lambda_{1,N},\ldots,\lambda_{N,N}\) be the eigenvalues of \(\Sigma_{X,N}^{1/2}\Sigma_{W,N}^{-1}\Sigma_{X,N}^{1/2}\).  Congruence in (6), together with equality of the spectra of \(\Sigma_{X,N}^{1/2}\Sigma_{W,N}^{-1}\Sigma_{X,N}^{1/2}\) and \(\Sigma_{W,N}^{-1/2}\Sigma_{X,N}\Sigma_{W,N}^{-1/2}\), gives
\[
 0<\lambda_{i,N}\leq C_\star,
 \qquad
 \sum_{i=1}^N\lambda_{i,N}
 =\operatorname{tr}(\Sigma_{W,N}^{-1}\Sigma_{X,N})
 \leq NC_\star.
 \tag{7}
\]

Choose \(\beta_0\in\operatorname{int}(\mathcal B)\), which is permitted by the slope-space assumption.  For a fixed \(q>0\), put
\[
 \Delta_N=qN^{-1/2},
 \qquad
 \beta_{1,N}=\beta_0+\Delta_N.
 \tag{8}
\]
Because \(\beta_0\) is an interior point, \(\beta_{1,N}\in\operatorname{int}(\mathcal B)\) for all sufficiently large \(N\).  Let \(P_{0,N}\) and \(P_{1,N}\) be the observed laws in \(\mathcal E_0\) with slopes \(\beta_0\) and \(\beta_{1,N}\), respectively, and with every other component fixed.  Although \(\beta_{1,N}\) changes with \(N\), at each \(N\) it is the \(N\)-th marginal of the class member whose slope is held equal to \(\beta_{1,N}\) at every sample size.  Thus both laws are among the alternatives over which the pointwise infimum and supremum in the theorem are taken.

Under \(P_{0,N}\), the vectors \(X_N\) and \(W_N=Y_N-\beta_0X_N\) are independent centered Gaussians.  Conditionally on \(X_N\), changing the slope by \(\Delta_N\) translates the mean of \(Y_N\) by \(\Delta_NX_N\) and leaves the conditional covariance equal to the invertible matrix \(\Sigma_{W,N}\).  Direct division of the two conditional Gaussian densities therefore gives the exact likelihood ratio on the complete observed vector \((X_N,Y_N)\):
\[
 L_N:=\frac{dP_{1,N}}{dP_{0,N}}
 =\exp\left\{\Delta_NX_N^{\mathsf T}\Sigma_{W,N}^{-1}W_N
 -\frac{\Delta_N^2}{2}X_N^{\mathsf T}\Sigma_{W,N}^{-1}X_N\right\}.
 \tag{9}
\]
Taking \(P_{0,N}\)-expectations of minus the logarithm and using independence and centering yields the exact divergence
\[
 \operatorname{KL}(P_{0,N},P_{1,N})
 =\frac{\Delta_N^2}{2}
 \operatorname{tr}(\Sigma_{W,N}^{-1}\Sigma_{X,N})
 \leq\frac{q^2C_\star}{2}.
 \tag{10}
\]
For the total-variation bound, conditioning the square of (9) on \(X_N\) and completing the one-dimensional Gaussian square gives
\[
 E_{0,N}[L_N^2\mid X_N]
 =\exp\{\Delta_N^2X_N^{\mathsf T}\Sigma_{W,N}^{-1}X_N\}.
\]
After diagonalizing the quadratic form against a standard normal vector, the elementary Gaussian integrals are finite when \(2\Delta_N^2C_\star<1\) and give
\[
 E_{0,N}L_N^2
 =\prod_{i=1}^N(1-2\Delta_N^2\lambda_{i,N})^{-1/2}.
 \tag{11}
\]
For all sufficiently large \(N\), \(2q^2C_\star/N\leq1/2\).  Since \(-\log(1-x)\leq2x\) for \(0\leq x\leq1/2\), (7), (8), and (11) imply
\[
 \log E_{0,N}L_N^2
 \leq\frac{2q^2}{N}\sum_{i=1}^N\lambda_{i,N}
 \leq2q^2C_\star.
 \tag{12}
\]
The definition of total variation and Cauchy--Schwarz now yield
\[
 \operatorname{TV}(P_{0,N},P_{1,N})
 =\frac12E_{0,N}|L_N-1|
 \leq\frac12\{E_{0,N}(L_N-1)^2\}^{1/2}
 \leq\frac12\{e^{2q^2C_\star}-1\}^{1/2}.
 \tag{13}
\]
Equations (9)--(13) concern the exact actual-design observed likelihood.  They use neither independent Fourier coordinates nor a boundary, aliasing, stencil, or Mat\'ern remainder approximation; those issues are absent from this converse subexperiment.

We first prove the confidence-set assertion.  By the nominal-level assumption,
\[
 \eta_\alpha=1-2\alpha>0.
\]
Choose
\[
 q_\alpha=left\{\frac{1}{2C_\star}
 \log\left(1+\frac{\eta_\alpha^2}{4}\right)\right\}^{1/2}>0.
 \tag{14}
\]
Equation (13) gives \(\operatorname{TV}(P_{0,N},P_{1,N})\leq\eta_\alpha/4\) eventually.  For an arbitrary candidate sequence of measurable random sets, define
\[
 a_N=\inf_{(P_N,\mathcal S_N)\in\mathcal M}P_N\{\beta\in I_N\}.
\]
The stated honesty condition is \(\liminf_Na_N\geq1-\alpha\).  Since both alternatives lie in \(\mathcal E_0\subset\mathcal M\), total variation and the union bound give
\[
 \begin{aligned}
 P_{0,N}\{\beta_0,\beta_{1,N}\in I_N\}
 &\geq P_{0,N}\{\beta_0\in I_N\}
       +P_{0,N}\{\beta_{1,N}\in I_N\}-1\\
 &\geq 2a_N-1-\operatorname{TV}(P_{0,N},P_{1,N}).
 \end{aligned}
 \tag{15}
\]
On this event, the definition of diameter and (8) give \(\operatorname{diam}(I_N)\geq\Delta_N\); no connectedness of \(I_N\) is required.  Therefore (14)--(15) imply
\[
 \begin{aligned}
 \liminf_{N\to\infty}\sqrt N
 \sup_{(P_N,\mathcal S_N)\in\mathcal M}
 E_{P_N}[\operatorname{diam}(I_N)]
 &\geq q_\alpha\liminf_{N\to\infty}
 P_{0,N}\{\beta_0,\beta_{1,N}\in I_N\}\\
 &\geq\frac34q_\alpha\eta_\alpha
 =:c_\alpha>0.
 \end{aligned}
 \tag{16}
\]

Finally, let \(\widetilde\beta_N\) be arbitrary and choose
\[
 q_0=\left(\frac{\log2}{2C_\star}\right)^{1/2}.
\]
Equation (13) gives \(\operatorname{TV}(P_{0,N},P_{1,N})\leq1/2\) eventually.  Define the test
\[
 \phi_N=\mathbf1\!\left\{\widetilde\beta_N\geq
 \frac{\beta_0+\beta_{1,N}}2\right\}.
\]
On \(\{\phi_N=1\}\) under \(P_{0,N}\), the squared error is at least \(\Delta_N^2/4\); on \(\{\phi_N=0\}\) under \(P_{1,N}\), it is at least the same amount.  Writing \(\beta_{0,N}=\beta_0\) and
\[
 R_{k,N}=E_{k,N}[(\widetilde\beta_N-\beta_{k,N})^2],
\]
we obtain
\[
 \begin{aligned}
 \max(R_{0,N},R_{1,N})
 &\geq\frac{\Delta_N^2}{8}
 \{P_{0,N}(\phi_N=1)+P_{1,N}(\phi_N=0)\}\\
 &\geq\frac{\Delta_N^2}{8}
 \{1-\operatorname{TV}(P_{0,N},P_{1,N})\}
 \geq\frac{q_0^2}{16N}.
 \end{aligned}
 \tag{17}
\]
The middle inequality is the testing identity \(P_0(\phi=1)+P_1(\phi=0)\geq1-\operatorname{TV}(P_0,P_1)\).  Since both alternatives belong to \(\mathcal M\), (17) yields
\[
 \liminf_{N\to\infty}N
 \sup_{(P_N,\mathcal S_N)\in\mathcal M}
 E_{P_N}[(\widetilde\beta_N-\beta)^2]
 \geq\frac{q_0^2}{16}
 =\frac{\log2}{32C_\star}=:c>0.
\]
Together with (16), this proves both assertions.

%%% FIELD statement-superclass
Let \(\mathfrak M^+\) be any statistical model of triangular arrays with the same observed vectors \((X(s_{i,N}),Y(s_{i,N}))_{i=1}^N\) and the same scalar slope target \(\beta\).  Suppose there exist \(d\in\{1,2\}\), a bounded Lipschitz domain \(D\), one deterministic design sequence satisfying \(c_mN^{-1/d}\leq h_N\leq C_mN^{-1/d}\) and \(\delta_N\geq c_sh_N\), one fixed tuple \((c_X,c_W,\alpha_X,\alpha_W,\alpha_C)\in\mathcal K_{\mathrm{nuis}}\), one \(p^\circ\in[p_{\min},p_{\max}]\), and numbers \(\beta^\circ\) and \(\varepsilon>0\) with \([\beta^\circ-\varepsilon,\beta^\circ+\varepsilon]\subset\operatorname{int}(\mathcal B)\), such that, for every \(b\) in this interval, \(\mathfrak M^+\) contains the full triangular-array law on that design generated by independent centered stationary Gaussian fields with marginal spectra \(f_X\) and \(f_W\), zero cross-spectrum \(A=0\), and response \(Y=bX+W\).  Then, for every \(0<\alpha<1/2\), there are constants \(c^+_\alpha,c^+>0\) such that every sequence \((I_N)_{N\geq1}\) of measurable random subsets of \(\mathbb R\), with measurable diameter and \(\operatorname{diam}(\varnothing)=0\), satisfying
\[
 \liminf_{N\to\infty}\inf_{(P_N,\mathcal S_N)\in\mathfrak M^+}
 P_N\{\beta\in I_N\}\geq1-\alpha,
\]
obeys
\[
 \liminf_{N\to\infty}\sqrt N
 \sup_{(P_N,\mathcal S_N)\in\mathfrak M^+}
 E_{P_N}[\operatorname{diam}(I_N)]\geq c^+_\alpha.
\]
Moreover, every estimator sequence \((\widetilde\beta_N)_{N\geq1}\) obeys
\[
 \liminf_{N\to\infty}N
 \sup_{(P_N,\mathcal S_N)\in\mathfrak M^+}
 E_{P_N}[(\widetilde\beta_N-\beta)^2]\geq c^+.
\]

%%% FIELD proof-superclass
Fix the embedded subexperiment asserted in the premise.  Its nuisance tuple is fixed and strictly positive, so the constant
\[
 C_\star=\frac{c_X}{c_W}
 \max\left\{1,\left(\frac{\alpha_W^2}{\alpha_X^2}\right)^{(d+1)/2}\right\}
\]
is finite and positive.  The exact actual-design likelihood calculation in the honest-confidence-set lower-bound theorem uses only this independent \(A=0\) subexperiment.  In particular, for \(\beta_{0,N}=\beta^\circ\), \(\beta_{1,N}=\beta^\circ+qN^{-1/2}\), and all sufficiently large \(N\) such that \(qN^{-1/2}<\varepsilon\), it gives
\[
 \operatorname{TV}(P_{0,N},P_{1,N})
 \leq\frac12\{e^{2q^2C_\star}-1\}^{1/2}.
 \tag{18}
\]
Both laws belong to \(\mathfrak M^+\) by the explicit inclusion premise.  This is the only point at which the ambient class enters the lower-bound proof.

Let \(a_N^+=\inf_{(P_N,\mathcal S_N)\in\mathfrak M^+}P_N\{\beta\in I_N\}\), put \(\eta_\alpha=1-2\alpha\), and choose
\[
 q_\alpha=\left\{\frac{1}{2C_\star}
 \log\left(1+\frac{\eta_\alpha^2}{4}\right)\right\}^{1/2}.
\]
Uniform honesty over \(\mathfrak M^+\) implies \(\liminf_Na_N^+\geq1-\alpha\).  The same total-variation and union-bound calculation, now using the two embedded laws, gives
\[
 P_{0,N}\{\beta_{0,N},\beta_{1,N}\in I_N\}
 \geq2a_N^+-1-\operatorname{TV}(P_{0,N},P_{1,N}).
\]
On this event \(\operatorname{diam}(I_N)\geq q_\alpha N^{-1/2}\).  Equation (18) therefore yields
\[
 \liminf_{N\to\infty}\sqrt N
 \sup_{(P_N,\mathcal S_N)\in\mathfrak M^+}
 E_{P_N}[\operatorname{diam}(I_N)]
 \geq\frac34q_\alpha(1-2\alpha)=:c_\alpha^+>0.
\]

For squared risk, take \(q_0=\{\log2/(2C_\star)\}^{1/2}\).  Then (18) is at most \(1/2\).  Thresholding any \(\widetilde\beta_N\) at \((\beta_{0,N}+\beta_{1,N})/2\) gives, exactly as in the theorem,
\[
 \sup_{(P_N,\mathcal S_N)\in\mathfrak M^+}
 E_{P_N}[(\widetilde\beta_N-\beta)^2]
 \geq\frac{q_0^2}{16N}
\]
for all sufficiently large \(N\).  Hence the second conclusion holds with
\[
 c^+=\frac{\log2}{32C_\star}>0.
\]
This is converse monotonicity with respect to inclusion of the exact subexperiment, not a claim that any named broader spatial model automatically contains it.
